\documentclass{article}
\usepackage{iclr2027_conference,times}
\input{math_commands.tex}

\usepackage{booktabs}
\usepackage{graphicx}
\usepackage{hyperref}
\usepackage{microtype}
\usepackage{url}
\usepackage{xcolor}

\title{LAD-Generic: Parameter-Efficient Conversion of\\Instruction-Tuned Language Models to Diffusion}

% Keep the submission anonymous. Do not uncomment \iclrfinalcopy until camera ready.
\author{Anonymous authors}

\newcommand{\method}{LAD-Generic}
\newcommand{\xxx}[1]{\textcolor{red}{\textbf{XXX: #1}}}
\newcommand{\na}{--}

\begin{document}
\maketitle

\begin{abstract}
Masked diffusion language models permit parallel token updates and flexible generation orders, but leading systems require costly pretraining or continued full-model training. We ask whether diffusion generation can instead be added to existing instruction-tuned autoregressive models with one generic parameter-efficient recipe. \method{} replaces causal attention with bidirectional attention and trains rank-1024 query/value LoRA adapters under a masked denoising objective, leaving 94.2--95.9\% of parameters frozen. Applied without architecture-specific tuning to Gemma~2~9B, Llama~3.1~8B, Qwen2.5~7B, and Ministral~8B, each 25,000-update conversion sees 400,000 examples, or at most 102.4M token positions, and runs on one GPU in a Google Colab environment. On our preliminary nine-task evaluation, Gemma is the strongest converted model with 56.4\% macro accuracy, compared with 55.1\% for LLaDA~8B Instruct in the same harness. Continued Llama training from 25k to 50k updates does not improve the nine-task macro average (51.3\% to 48.9\%), but preliminarily improves the parallel decoding regime, where the number of denoising network evaluations (NFE) is smaller than the 128-token output length: 32-step generation perplexity falls from 30.05 to 16.80 and 64-step perplexity from 10.50 to 8.71. Across all four families, the converted models generate coherent answers from fully masked sequences, demonstrating an accessible alternative to diffusion pretraining from scratch.
\end{abstract}

\section{Introduction}

Autoregressive (AR) language models generate a sequence from left to right. This method has enabled effective scaling, but it fixes the generation order, limiting the model to unidirectional attention. Masked diffusion language models instead begin with a corrupted sequence and iteratively reconstruct multiple positions using bidirectional context \citep{austin2021structured,sahoo2024simple,shi2024simplified}. Their number of network evaluations can be chosen independently of output length, and early erroneous predictions can be revised. However, training or adapting models for diffusive generation has thus far been a computationally intensive and expensive procedure.

The authors of LLaDA were the first to show that masked diffusion can scale to an 8B-parameter general-purpose language model \citep{nie2025llada}. While they based the architecture on the Llama 3.1 8B model, they pre-trained it from scratch using 2.3T tokens and 130,000 GPU-hours. DiffuLLaMA demonstrated a more efficient route, converting a similarly sized pre-trained AR model through continued diffusion training \citep{gong2025diffullama}. Since then, various methods have been proposed to improve efficiency and task retention. Dream explicitly builds on DiffuLLaMA's AR initialization and adds context-adaptive token-level noise rescheduling \citep{ye2025dream}; Efficient-DLM combines blockwise attention with position-dependent masking \citep{fu2025efficientdlm}; and iLLaDA returns to from-scratch diffusion pretraining at substantially larger scale \citep{nie2026illada}. LAD was the first to show that low-rank adaptation (LoRA) can induce diffusion generation with substantially fewer trainable parameters \citep{kuiper2025lad,hu2022lora}. It remains unclear how reliably this minimal recipe transfers across architecture families and how much capability is retained relative to each model's own AR parent.

Our aim is not merely to avoid updating every parameter, but to provide a practical alternative to training a separate diffusion foundation model for every model family. Reusing existing AR knowledge reduces the incremental data, accelerator count, and financial barrier to experimentation. This emphasis follows calls to treat computational efficiency and access as research outcomes alongside accuracy \citep{strubell2019energy}. It also motivates a sustainability benefit, although compute volume alone is not a carbon measurement: emissions depend strongly on hardware, datacenter efficiency, location, and energy supply \citep{patterson2021carbon}. We therefore report bounded incremental training scale and hardware while reserving quantitative environmental claims for future energy measurements.

We study this question with \method{}, a generalizable conversion and evaluation pipeline for four instruction-tuned decoder families: Gemma~2~9B, Llama~3.1~8B, Qwen2.5~7B, and Ministral~8B. At the time the experimental design was frozen, we selected the latest openly available text-only instruction checkpoint from each chosen family at a comparable nominal scale that fit the same single-GPU pipeline. Excluding multimodal variants avoids vision components and modality-specific training as additional confounders. Every backbone receives the same bidirectional attention intervention, LoRA target modules, and masked-diffusion objective. This paired design separates conversion quality from the strength of the starting checkpoint. We directly evaluate LLaDA~8B Instruct because its standard masked-diffusion objective, instruction-tuning setup, and overlapping benchmark suite make it the cleanest external control. We additionally report published DiffuLLaMA, Dream-Instruct, iLLaDA-Instruct, and Efficient-DLM values, but keep them protocol-separated: their additional attention, masking, scheduling, data, and inference techniques would obscure rather than strengthen a controlled comparison.

Our contributions are:
\begin{itemize}
    \item a single, architecture-agnostic recipe that converts four instruction-tuned AR model families to masked diffusion using only Q/V LoRA parameters on one Google Colab GPU;
    \item a paired evaluation of each diffusion model and its frozen AR parent across general knowledge, reasoning, mathematics, and code benchmarks;
    \item a controlled analysis of base-model choice, 25k versus 50k updates, denoising-step count, validation loss, and externally scored generation perplexity; and
    \item a transparent comparison between results produced by our common harness and the published values reported with the LLaDA evaluation protocol.
\end{itemize}

\section{Related work}

Discrete diffusion models corrupt and reconstruct categorical sequences rather than adding Gaussian noise \citep{austin2021structured}. Recent masked formulations simplify this process to absorbing-state masking and yield tractable likelihood-based objectives \citep{sahoo2024simple,shi2024simplified}. LLaDA scaled this formulation to language-model pretraining and instruction following \citep{nie2025llada}; iLLaDA retains full bidirectionality but scales pretraining to 12T tokens and uses a 25B-token instruction corpus for 12 epochs \citep{nie2026illada}.

AR-to-diffusion conversion amortizes the knowledge already acquired by an AR checkpoint. DiffuLLaMA established this route through full-model continued pretraining, using 65B tokens at 7B scale \citep{gong2025diffullama}. Dream does not independently introduce AR initialization: it explicitly adopts DiffuLLaMA's shifted-prediction initialization, continues Qwen2.5 for 580B tokens, adds context-adaptive token-level noise rescheduling, and then tunes on 1.8M instruction pairs for three epochs \citep{ye2025dream}. Efficient-DLM preserves causality across blocks, exposes clean preceding blocks, and biases masking toward later positions; its 8B model receives 500B tokens of continued training \citep{fu2025efficientdlm}. These methods target stronger task retention and/or faster decoding, whereas \method{} deliberately tests a smaller intervention: fully bidirectional attention and Q/V-only LoRA on instruction data.

Hybrid and blockwise models offer a different compromise between bidirectionality and deployability. Block Diffusion is autoregressive across blocks but diffusive within each block, enabling variable-length generation and inter-block KV caching \citep{arriola2025block}. Fast-dLLM v2 converts Qwen2.5 models with approximately 1B fine-tuning tokens and combines blockwise complementary masking with hierarchical caching \citep{wu2025fastdllmv2}. Sequential Diffusion Language Models instead learn next-sequence prediction and dynamically finalize variable-length subsequences, reporting conversion with 3.5M training samples \citep{liu2025sdlm}. These systems retain more left-to-right structure than our fully bidirectional answer block, making them especially relevant to speed, caching, and arbitrary-length generation but less direct as controls for our intervention.

Several recent systems clarify adjacent goals. Uno freezes an AR model and learns lightweight LoRA diffusion-draft weights by consistency distillation, then uses the unchanged AR model as a verifier; its $\Psi$-Spec sampler therefore provides lossless acceleration of the AR distribution rather than replacing it with a standalone diffusion distribution \citep{sahoo2026uno}. The dLLM framework standardizes training, inference, and evaluation across LLaDA- and Dream-style models. Importantly, its rank-128 LoRA experiment is reasoning SFT of these already-diffusive checkpoints, not LoRA-based AR-to-diffusion conversion. Its separate Tiny-A2D experiment converts Qwen3-0.6B with instruction SFT on 64 A100 GPUs and does not report LoRA as the conversion mechanism \citep{zhou2026dllm}. NaRA likewise uses noise-conditioned LoRA to adapt existing diffusion LLMs rather than to induce diffusion behavior in an AR checkpoint \citep{wang2026nara}. These distinctions isolate our central question: can low-rank updates themselves perform the large-model paradigm conversion?

Data efficiency is not only a conversion issue. Controlled pretraining experiments find that masked diffusion can benefit unusually strongly from repeated passes over limited unique data, partly because independently resampled masks provide implicit augmentation \citep{ni2025superdata}. This motivates checkpoint-wise and multi-epoch study, but does not guarantee monotonic gains under our instruction-only conversion setting.

Our focus is complementary: we hold the adaptation mechanism as constant as the architectures permit, repeat it across four families, and benchmark each converted model against the exact AR checkpoint from which it originated. This paired comparison is essential because raw diffusion scores otherwise conflate conversion quality with the prior capability of the parent model.

Table~\ref{tab:compute} places the incremental training budgets side by side. Token and example counts are more defensible than cross-generation GPU-hour normalization, so we retain NR where papers do not report elapsed time. For conversion methods, these budgets exclude the substantial earlier cost of training the AR parent; for from-scratch methods they include diffusion pretraining. Consequently, the table compares the additional cost required to obtain diffusion behavior, not total lifecycle compute.

\begin{table*}[t]
\caption{Reported training scale for relevant large diffusion systems. ``Full'' means the main diffusion model weights are updated; LoRA denotes parameter-efficient diffusion weights. Hardware generations and sequence lengths differ, so GPU-hours are not interchangeable. NR: not reported.}
\label{tab:compute}
\centering
\scriptsize
\resizebox{\textwidth}{!}{%
\begin{tabular}{llllll}
\toprule
System & Starting point & Diffusion-stage data & Updated weights & Hardware & Reported time \\
\midrule
\method{} 25k (ours) & instruction AR & $\leq$102.4M positions & Q/V LoRA & 1$\times$RTX Pro 6000 (Colab) & \xxx{hours} \\
\method{} 50k (ours) & instruction AR & $\leq$204.8M positions & Q/V LoRA & 1$\times$RTX Pro 6000 (Colab) & \xxx{hours} \\
DiffuLLaMA 7B & LLaMA2 7B & 65B tokens & Full & 64$\times$GH200 & NR \\
LLaDA 8B & random init. & 2.3T tokens + 4.5M SFT pairs & Full & H800 & 130,000 GPU-h (pretrain) \\
Dream 7B & Qwen2.5 7B & 580B tokens + 1.8M pairs $\times3$ & Full & NR & NR \\
iLLaDA 8B & random init. & 12T tokens + 25B-token SFT $\times12$ & Full & NR & NR \\
Efficient-DLM 8B & Qwen3 8B & 500B tokens & Full & 128$\times$H100 & NR \\
Uno 8B & frozen AR & 7B distillation tokens & LoRA diffusion weights & 64$\times$H200 & 3,840 GPU-h \\
dLLM reasoning SFT & LLaDA/Dream (DLM) & s1K $\times20$ epochs & rank-128 LoRA & 8$\times$A100 & NR \\
dLLM Tiny-A2D & Qwen3 0.6B AR & instruction mixture $\times10$ epochs & NR (not LoRA-reported) & 64$\times$A100 & NR \\
Fast-dLLM v2 7B & Qwen2.5 7B AR & $\sim$1B tokens & Full & 64$\times$A100 & 768 GPU-h \\
SDLM & instruction AR & 3.5M samples & NR & NR & NR \\
\bottomrule
\end{tabular}}
\end{table*}

\section{Method}

\subsection{Conditional masked denoising}

Let $p$ be an observed instruction prompt and $x_0=(x_0^1,\ldots,x_0^L)$ its clean response region. For each example, we sample
$t\sim\mathcal{U}[10^{-3},1]$ and independently replace each eligible response token with a single tokenizer-native mask representation with probability $t$. Denote the resulting sequence by $x_t$ and the sampled mask positions by $M_t$. The model is trained with same-position cross-entropy,
\begin{equation}
    \mathcal{L}(\theta) =
    -\mathbb{E}_{t,x_0,x_t}
    \left[
    \frac{1}{t|E|}\sum_{i\in M_t}
    \log p_{\phi,\theta}(x_0^i\mid p,x_t)
    \right],
    \label{eq:objective}
\end{equation}
where $E$ is the eligible answer-and-padding region, $\phi$ denotes the frozen parent parameters, and $\theta$ the LoRA parameters. Prompt and chat-template tokens are always visible and never enter the loss. Training corruption is resampled online; validation corruption is deterministic from the example index and seed.

The clean terminating EOS is included as a genuine answer token. To teach the fixed-width generator to emit termination beyond the answer, dynamic batch padding is represented by repeated EOS tokens; these positions are also eligible for masking and loss. Thus padding is part of the learned denoising target rather than conventional ignored padding. This design choice should be interpreted separately from the semantic answer-token objective and is ablated in \xxx{appendix/table after final ablation selection}.

\subsection{Bidirectional conversion and trainable parameters}

We replace the causal mask with an all-zero four-dimensional additive attention mask, so every prompt, answer, and EOS-padding position can attend to every position. The architecture and output vocabulary remain unchanged. We attach LoRA adapters to every self-attention query and value projection, with rank $r=1024$, scaling $\alpha=1024$, and dropout zero. Backbone weights, token embeddings, the language-model head, biases, and normalization layers are frozen. Each model reuses an existing single-token string from its tokenizer (\texttt{MASK} for Gemma, Llama, and Qwen; \texttt{<?>} for Ministral), avoiding vocabulary expansion. At experiment-design time, these were the latest openly available text-only instruction checkpoints from the four selected families at a comparable nominal scale; using text-only variants keeps the conversion interface uniform and avoids multimodal components as a confounder.

Table~\ref{tab:models} reports the parameter audits saved by the training runs. The trainable fraction varies because the four attention implementations have different projection shapes and layer counts.

\begin{table}[t]
\caption{Converted parent checkpoints and audited trainable parameter counts. All trainable parameters are Q/V LoRA weights.}
\label{tab:models}
\centering
\small
\begin{tabular}{lrrr}
\toprule
Parent checkpoint & Total parameters & Trainable & Trainable (\%) \\
\midrule
Gemma 2 9B IT & 10.732B & 572.5M & 5.33 \\
Llama 3.1 8B Instruct & 8.466B & 436.2M & 5.15 \\
Qwen2.5 7B Instruct & 7.939B & 323.0M & 4.07 \\
Ministral 8B Instruct 2410 & 8.511B & 490.7M & 5.77 \\
\bottomrule
\end{tabular}
\end{table}

After training, we fold each LoRA branch into its parent projection, $W_{\mathrm{merged}}=W+(\alpha/r)BA$, and remove the adapter modules. The resulting standalone checkpoint retains bidirectional inference but has the base architecture's parameter count: Llama decreases from 8.466B adapter-loaded parameters to 8.030B merged parameters (16.06~GB in BF16). Thus the model-size increase is temporary, although optimizer states and activations still add training-time memory.

\subsection{Training data}

All main models use \texttt{Ruurd/LAD-training-1m-256}. The training pipeline shuffles the source with seed 42 and consumes 400,000 usable examples per 25,000-update stage. Examples are retokenized with each parent model's native chat template and tokenizer. Sequences are dynamically padded to a multiple of 16 and limited to 256 tokens; only the current answer region is corrupted. The dataset builder groups sources into general instruction following, multiple-choice reasoning, mathematics, and code, deduplicates normalized prompts, and blocks exact normalized matches with represented benchmark validation/test prompts. \xxx{Archive the exact Hub revision and manifest used by these original runs, then insert the final category/source counts and long-answer filtering policy.}

The four 25k models see the same seed-selected 400k-example prefix. For the 50k Llama condition, the 25k final adapter initializes a second run over the next 400,000 examples from the same stable dataset shuffle. The second stage resamples corruption masks and restarts AdamW and its learning-rate schedule; it is therefore continued adapter training, not a restoration of optimizer state.

\subsection{Optimization}

Every run uses batch size 16 with no gradient accumulation, AdamW, learning rate $10^{-5}$, weight decay 0.01, gradient clipping at 0.5, 100 warmup updates, and cosine decay. The maximum is 25,000 optimizer updates per stage. Training uses BF16 master weights and Transformer Engine HYBRID FP8 execution; evaluation runs in BF16. All reported conversions were executed in a single-GPU Google Colab runtime using one NVIDIA RTX PRO 6000 Blackwell Server Edition, with no multi-GPU or multi-node parallelism. The code, memory strategy, sequence length, and adapter scope were explicitly designed around this constraint. We save model checkpoints every 1,000 updates (5,000 in the continued stage), evaluate denoising loss every 250 updates on 200 deterministic validation examples, and use seed 42. \xxx{Insert wall-clock duration, GPU-hours, peak memory, Colab tier, and software versions.}

\subsection{Diffusion generation}

Generation begins from a fully masked answer block. At each reverse step the model predicts every masked position, after which the least-confident positions are retained for reconsideration according to the LLaDA fixed-transfer schedule; newly accepted positions become permanent context. For the primary benchmark evaluation we use greedy token selection, low-confidence remasking, and task-specific output/step budgets (Table~\ref{tab:budgets}). Open-ended evaluation generates 128 tokens with 64 denoising steps, temperature 0.7, and delayed retention of EOS/EOT predictions. \xxx{Confirm these settings against the immutable benchmark run manifests associated with every workbook column.}

\begin{table}[t]
\caption{Current task budgets for the shared LLaDA-style diffusion evaluator. These must be verified against the final saved benchmark manifests.}
\label{tab:budgets}
\centering
\small
\resizebox{\columnwidth}{!}{%
\begin{tabular}{lrrrrrrrrr}
\toprule
Task & MMLU & MMLU-P & HSwag & ARC-C & GSM8K & MATH & GPQA & HEval & MBPP \\
\midrule
Tokens/steps & 64 & 128 & 64 & 64 & 256 & 256 & 128 & 256 & 128 \\
\bottomrule
\end{tabular}}
\end{table}

\section{Evaluation}

\subsection{Paired models and baselines}

For Gemma, Llama, Qwen, and Ministral we compare the diffusion adapter with greedy generation from its unchanged AR parent. The Llama parent is shared by the 25k and 50k diffusion checkpoints. LLaDA~8B Instruct is evaluated as an external diffusion baseline. We report two explicitly separated LLaDA comparisons: (i) LLaDA run through our prompts, sampling code, subsets, and scorers; and (ii) values transcribed from the LLaDA paper/evaluation suite. The latter are contextual references rather than directly matched measurements.

\subsection{Benchmarks and scoring}

Our suite contains MMLU, MMLU-Pro, HellaSwag, ARC-Challenge, GSM8K, MATH-500, GPQA-Main, HumanEval, and sanitized MBPP \citep{hendrycks2021mmlu,wang2024mmlupro,zellers2019hellaswag,clark2018arc,cobbe2021gsm8k,hendrycks2021math,rein2023gpqa,chen2021codex,austin2021mbpp}. HellaSwag uses its labelled validation split, ARC-Challenge and MATH-500 use test, and GPQA-Main uses the release's train-named evaluation split. Multiple-choice answers are extracted by label; GSM8K uses the final numeric answer; MATH uses normalized boxed answers; and HumanEval/MBPP execute tests in isolated local subprocesses. The current workbook records 50 examples per task unless marked $\dagger$, where 250 were evaluated. This is a preliminary subset evaluation; the final paper will report \xxx{full-set results or a single predeclared sampling design, with uncertainty intervals}.

For 30 fixed open-ended prompts, Microsoft Phi-4 scores each generated answer autoregressively. Our primary aggregate is token-weighted perplexity, $\exp(\sum_j \mathrm{NLL}_j/\sum_j N_j)$, so each scored token contributes equally; per-answer means and medians are retained as supplementary diagnostics. These are external-reference fluency measures, not likelihoods of the diffusion model itself. This distinction is important: generative perplexity inherits the reference model's preferences and does not penalize loss of sample diversity \citep{turok2026duel}. We therefore also compute sliding model-token Distinct-1/2/3 after excluding special tokens. Higher Distinct-$n$ means a larger fraction of unique overlapping $n$-grams. \xxx{Recompute or replace workbook values labelled ``1-gram repetition,'' which predate the final Distinct-$n$ definition.}

\section{Results}

\subsection{Preliminary benchmark accuracy}

Table~\ref{tab:accuracy} transcribes the current \texttt{Results.xlsx} workbook. Empty cells are deliberately retained as \texttt{XXX}: in particular, several paired AR evaluations are still running. These preliminary values should not yet be used for significance claims.

\begin{table*}[t]
\caption{Preliminary exact-match/pass@1 accuracy in our matched harness. D: converted diffusion model; AR: unchanged parent. Llama-50k shares the Llama-25k AR parent. Bold marks the best diffusion result in each row; underline marks the best converted \method{} result (ties retained). Thus a bold underlined value is best both overall and among our conversions. $\dagger$ denotes 250 examples and other populated cells use 50. Parenthesized averages cover only the indicated number of populated tasks; XXX is unfinished.}
\label{tab:accuracy}
\centering
\scriptsize
\resizebox{\textwidth}{!}{%
\begin{tabular}{lccccccccccc}
\toprule
& \multicolumn{2}{c}{Gemma 25k} & \multicolumn{3}{c}{Llama} & \multicolumn{2}{c}{Qwen 25k} & \multicolumn{2}{c}{Ministral 25k} & \multicolumn{2}{c}{LLaDA} \\
Task & D & AR & D-25k & D-50k & AR & D & AR & D & AR & Ours & Published \\
\midrule
ARC-C      & \textbf{\underline{.880}} & XXX & .828$^\dagger$ & .780 & .776$^\dagger$ & .840 & XXX & .820 & .740 & \textbf{.880} & .885 \\
GPQA       & .360 & XXX & .280 & .220 & XXX & \textbf{\underline{.380}} & XXX & .240 & XXX & .300 & .333 \\
GSM8K      & .600 & XXX & \textbf{\underline{.652}}$^\dagger$ & .620 & .856$^\dagger$ & .560 & XXX & .600 & .860 & .640 & .694 \\
HellaSwag  & .760 & XXX & .732$^\dagger$ & \textbf{\underline{.840}} & .588$^\dagger$ & .740 & XXX & .820 & .720 & .720 & .746 \\
HumanEval  & \underline{.500} & XXX & \underline{.500} & .380 & XXX & .320 & XXX & XXX & XXX & \textbf{.580} & .494 \\
MATH-500   & .300 & XXX & .220 & .180 & XXX & \underline{.360} & XXX & XXX & XXX & \textbf{.380} & .319 \\
MBPP       & \textbf{\underline{.600}} & XXX & .460 & .400 & XXX & .540 & XXX & XXX & XXX & .480 & .410 \\
MMLU       & \textbf{\underline{.640}} & .720 & .612$^\dagger$ & .620 & .560$^\dagger$ & .600 & .660 & .600 & .520 & .560 & .655 \\
MMLU-Pro   & \textbf{\underline{.440}} & .460 & .332$^\dagger$ & .360 & .448$^\dagger$ & .280 & .400 & .360 & .320 & .420 & .370 \\
\midrule
Average    & \textbf{\underline{.564}} & .590 (2) & .513 & .489 & .646 (5) & .513 & .530 (2) & .640 (5) & .632 (5) & .551 & .545 \\
\bottomrule
\end{tabular}}
\end{table*}

The provisional complete diffusion macro averages are 56.4\% for Gemma, 51.3\% for Llama-25k, 48.9\% for Llama-50k, 51.3\% for Qwen, and 55.1\% for LLaDA in our harness. Ministral's 64.0\% and the AR averages are explicitly partial. These means mix 50- and 250-example estimates and are descriptive only. Notably, the second Llama stage does not uniformly improve downstream accuracy despite reducing its 64-step generation perplexity (Section~\ref{sec:perplexity}), motivating checkpoint-wise evaluation rather than selecting solely by denoising loss or fluency.

\begin{table}[t]
\caption{Published diffusion-model accuracies transcribed from \texttt{Results.xlsx}; NR means not reported in the workbook. These values use the respective papers' protocols and are not directly comparable to Table~\ref{tab:accuracy}. Bold is the best reported value in a row; underline is the best excluding LLaDA. The final row reproduces each workbook column's stated average, whose task coverage differs across papers.}
\label{tab:published}
\centering
\scriptsize
\resizebox{\columnwidth}{!}{%
\begin{tabular}{lccccc}
\toprule
Task & LLaDA & DiffuLLaMA & Dream-Inst. & iLLaDA-Inst. & Efficient-DLM 8B \\
\midrule
ARC-C     & \textbf{.885} & NR & NR & NR & \underline{.619} \\
GPQA      & \textbf{.333} & NR & \underline{.330} & NR & NR \\
GSM8K     & .694 & NR & .810 & \textbf{\underline{.890}} & .883 \\
HellaSwag & \textbf{.746} & NR & NR & NR & \underline{.725} \\
HumanEval & .494 & NR & .555 & .659 & \textbf{\underline{.683}} \\
MATH      & .319 & NR & .392 & \textbf{\underline{.567}} & .501 \\
MBPP      & .410 & NR & .588 & .580 & \textbf{\underline{.610}} \\
MMLU      & .655 & NR & .670 & .716 & \textbf{\underline{.772}} \\
MMLU-Pro  & .370 & NR & .433 & \textbf{\underline{.523}} & NR \\
\midrule
Workbook average & .545 & NR & .495 & .609 & \textbf{\underline{.652}} \\
\bottomrule
\end{tabular}}
\end{table}

DiffuLLaMA is included for methodological completeness, but the workbook contains no compatible reported values for these nine instruction benchmarks, so its cells remain NR rather than importing scores from a different task suite. The published averages also have different denominators: the workbook averages all nine LLaDA tasks but excludes GSM8K, as well as unreported cells, for Dream, iLLaDA, and Efficient-DLM. They therefore summarize each column but must not be read as a controlled ranking.

\subsection{Open-ended fluency and denoising budget}
\label{sec:perplexity}

\begin{table}[t]
\caption{Token-weighted Phi-4 perplexity on 128-token open-ended generations at 64 denoising steps. Lower is better, but this is a secondary fluency diagnostic.}
\label{tab:ppl}
\centering
\small
\begin{tabular}{lrr}
\toprule
Model family/checkpoint & Diffusion & AR parent \\
\midrule
Gemma 2 9B, 25k & 8.017 & 3.011 \\
Llama 3.1 8B, 25k & 10.503 & 2.489 \\
Llama 3.1 8B, 50k & 8.709 & 2.489 \\
Qwen2.5 7B, 25k & 10.272 & 2.230 \\
Ministral 8B, 25k & 9.424 & 2.305 \\
LLaDA 8B Instruct & 4.065 & \na \\
\bottomrule
\end{tabular}
\end{table}

Increasing the denoising budget strongly improves this diagnostic. For example, Llama-25k changes from 30.049 at 32 steps to 10.503 at 64 and 3.753 at 128; Gemma-25k changes from 21.795 to 8.017 and 3.317. We call settings with fewer denoising network evaluations than generated tokens ($\mathrm{NFE}<L$) the \emph{parallel decoding regime}; unlike AR decoding, they permit more than one token to be finalized per model evaluation. Perplexity values are meaningful only under matched output length, sampler, temperature, and denoising budget.

We also benchmarked the saved Llama adapter before and after merging under identical 30-prompt, 128-token, 32-step settings (Appendix~\ref{app:merge}). Merging removes 436.2M deployed parameters while changing token-weighted perplexity from 19.381 to 19.283 ($-0.51\%$); Distinct-1/2/3 change from .495/.788/.899 to .498/.807/.918. This matched test provides no evidence of deterioration from merging.

\subsection{Progression with training updates}

Table~\ref{tab:progression} shows checkpoint-wise Phi-4 perplexity from the workbook. Improvements are not monotonic at every budget. The matched final-checkpoint comparison suggests that the second Llama stage primarily helps in the parallel decoding regime: for 128-token outputs, 25k to 50k training lowers perplexity from 30.049 to 16.797 at 32 NFE and from 10.503 to 8.709 at 64 NFE, while the 128-NFE value worsens from 3.753 to 4.876. Meanwhile, the preliminary factual-task macro average changes from 51.3\% to 48.9\%. Continued training therefore appears to improve low-budget parallel decoding rather than general factual capability; ongoing evaluations will test whether this pattern holds across models and seeds. Gemma reaches most of its 64-step improvement within 15--20k updates and slightly degrades by 25k. Final benchmark results will report fixed checkpoints rather than choosing a separate best checkpoint per task.

\begin{table}[t]
\caption{Training progression: token-weighted Phi-4 perplexity for 128-token generations.}
\label{tab:progression}
\centering
\scriptsize
\begin{tabular}{llrrrrrr}
\toprule
Model & Steps & 1k & 10k & 20k & 30k & 40k & 50k \\
\midrule
Llama & 32  & 15.502 & 23.138 & 17.091 & 14.737 & 17.201 & 16.797 \\
Llama & 64  & 10.256 & 10.015 & 9.646 & 9.839 & 8.716 & 8.382 \\
Llama & 128 & 5.795 & 5.287 & 4.763 & 4.664 & 4.420 & 4.876 \\
\midrule
Model & Steps & 1k & 5k & 10k & 15k & 20k & 25k \\
\midrule
Gemma & 32  & 19.206 & 13.235 & 11.547 & 12.413 & 10.398 & 11.688 \\
Gemma & 64  & 10.387 & 7.506 & 7.941 & 6.746 & 6.751 & 7.826 \\
Gemma & 128 & 5.852 & 4.457 & 4.313 & 3.454 & 3.941 & 3.872 \\
\bottomrule
\end{tabular}
\end{table}

\subsection{Base-model dependence and LLaDA comparison}

All four converted families produce usable diffusion generations, but the resulting quality is not architecture invariant. Gemma currently has the strongest complete diffusion macro average, while Qwen and Llama are similar on average but differ by task. The incomplete Ministral row already illustrates another important point: its strong multiple-choice and GSM8K scores do not establish code or MATH capability until those evaluations finish. \xxx{Replace this paragraph with paired D-minus-AR retention gaps and confidence intervals once every AR cell is populated.}

LLaDA scores 55.1\% on the nine-task unweighted average in our harness, compared with 54.5\% from the values transcribed as published. Agreement varies by task and should not be interpreted as an exact reproduction because prompts, subsets, generation lengths, answer extraction, and code execution may differ. The final submission will present our common-harness comparison as primary and Table~\ref{tab:published} only as external calibration. We chose LLaDA for direct external testing because it uses the uncomplicated response-masking objective closest to ours and reports the largest overlap with our task suite. Dream, iLLaDA, and Efficient-DLM add different data scales, noise weighting, SFT formats, scoring rules, attention patterns, or positional masking. Importing those techniques into the primary experiment would prevent us from attributing outcomes to parameter-efficient conversion itself.

\section{Discussion}

The central empirical finding is qualitative but robust across the selected runs: high-rank attention-only LoRA is sufficient to turn four independently developed causal instruction models into generators that reconstruct answers from a fully masked block. This supports the view that much of the knowledge needed for diffusion generation is already present in AR representations and vocabulary heads. Because each conversion fits on one Google Colab GPU rather than requiring a dedicated multi-accelerator cluster, the method also lowers the infrastructure threshold for reproducing, modifying, and teaching large-scale diffusion-language-model conversion. A capable paid cloud GPU is still required, so this is broader access rather than zero-cost access.

The merge result also shows that LoRA's parameter-count overhead is not permanent: the deployed Llama returns from 8.466B to 8.030B parameters without measurable degradation in this small matched test. This is not model compression---the full base-sized checkpoint remains---and more prompts and backbones are needed to bound small merge effects.

Conversion is not free, however. First, 323--573M trainable parameters are parameter-efficient relative to full-model training but substantially larger than conventional low-rank task adapters. Second, our secondary perplexity remains worse than each AR parent and than LLaDA at 64 steps. Third, additional Llama training does not improve the provisional downstream benchmark aggregate, but its largest fluency gains occur when $\mathrm{NFE}<L$. This suggests a more specific benefit: additional denoising training may teach the model to resolve multiple mutually dependent positions per evaluation, improving the quality--compute trade-off without adding corresponding benchmark capability. The 128-NFE result cautions that this interpretation remains preliminary. Finally, EOS-padding supervision makes the fixed-width objective different from a pure answer-only likelihood and deserves explicit ablation.

Perplexity alone is insufficient for selecting a diffusion generator. Our soft-frontier masking ablation lowered Phi-4 generative perplexity but substantially reduced sliding model-token Distinct-1/2/3 and produced visibly more repetitive answers (Appendix~\ref{app:frontier}). DiffuLLaMA similarly pairs external-model perplexity with distinct bigrams to check that low perplexity was not obtained through repetition \citep{gong2025diffullama}. More fundamentally, external generative perplexity reflects a separate AR model's preferences and ignores diversity, so repeated high-probability text can score deceptively well \citep{turok2026duel}. Theory likewise distinguishes token-level perplexity from whole-sequence correctness: a masked diffusion model may approach favorable perplexity with few steps while still requiring a step count that scales with sequence length for low sequence error \citep{feng2025benefit}. We therefore triangulate external perplexity with Distinct-1/2/3, task accuracy, and generated samples. Distinct-$n$ is itself only a repetition diagnostic, so grammaticality, coherence, and actual usability ultimately require systematic qualitative inspection rather than replacement by any single automatic score.

The literature comparison sharpens the trade-off. DiffuLLaMA and Dream reuse AR initialization but update the full model for 65B and 580B tokens, respectively; Efficient-DLM invests 500B tokens and architectural training choices specifically designed to retain causal-model behavior and enable cached block decoding \citep{gong2025diffullama,ye2025dream,fu2025efficientdlm}. iLLaDA's strong reported results come from 12T-token diffusion pretraining plus repeated large-scale SFT, a fundamentally different compute regime \citep{nie2026illada}. Fast-dLLM v2 and SDLM show that retaining blockwise or sequential structure can reduce conversion cost while enabling practical caching and variable-length generation \citep{wu2025fastdllmv2,liu2025sdlm}; our complementary question is how far a fully bidirectional model can be converted with adapter-only updates. Uno is closer to our parameter-efficient philosophy because it also trains LoRA diffusion weights, but its goal is different: an AR verifier guarantees samples from the original AR distribution, whereas \method{} seeks an independently usable fully bidirectional generator \citep{sahoo2026uno}. Thus the higher published scores in Table~\ref{tab:published} do not isolate a single superior masking objective; they bundle substantially different data budgets and modeling or decoding machinery.

The incremental scale is central to the intended contribution. A 25k \method{} run processes at most 102.4M token positions, versus tens or hundreds of billions of tokens for the full-model conversion systems in Table~\ref{tab:compute}, while leaving most parent parameters frozen. This comparison does not erase the sunk cost of AR pretraining, nor prove a proportional reduction in energy or emissions. It does show that an existing foundation checkpoint can be repurposed without repeating that pretraining or reserving a large cluster. That distinction matters both economically and for widening participation in diffusion-model research; direct power and carbon measurements remain necessary for an environmental comparison \citep{strubell2019energy,patterson2021carbon}.

This is also why we do not claim that Table~\ref{tab:published} is a leaderboard. Our direct LLaDA run holds prompts, subsets, extraction, and generation code fixed. The remaining columns are literature context, useful for locating the attainable performance range but unsuitable for causal comparison. A future study can add these systems to the shared harness one at a time, provided their task-specific inference procedures are faithfully reproduced---a requirement emphasized by the dLLM framework \citep{zhou2026dllm}.

\section{Limitations}

The current workbook combines 50- and 250-example benchmark subsets and does not yet contain every AR baseline. Unless expanded, these estimates have high sampling uncertainty. Only one training seed is available per model, so architecture and random-seed effects are not separated. The four backbones differ in tokenizer, pretraining data, alignment, attention implementation, and nominal parameter count. Although chosen as the latest comparable text-only releases when the study was designed, they are not guaranteed to remain the latest checkpoints at publication time. Phi-4 perplexity favors text that its own AR distribution considers probable, ignores diversity, and is not a diffusion likelihood. Distinct-$n$ detects local repetition but not factuality, grammar, or global coherence. Published LLaDA numbers are not directly comparable to our harness. We also have not yet established wall-clock speedups: $\mathrm{NFE}<L$ reduces sequential model calls but does not automatically imply lower latency because each call processes the full answer block. Finally, we did not instrument the runs for energy consumption or carbon intensity. Single-GPU execution and lower incremental training volume support accessibility and compute-efficiency claims, not a quantified sustainability claim.

\section{Conclusion}

\method{} provides a common, parameter-efficient route from instruction-tuned AR checkpoints to masked diffusion generators across Gemma, Llama, Qwen, and Ministral. The preliminary evidence shows broad transfer of the conversion mechanism, but also substantial dependence on the parent model, inference budget, and evaluation metric. By fitting every conversion on a single Google Colab GPU, the work offers a reproducible pathway for researchers who cannot retrain a foundation model or provision a large accelerator cluster. Continued training appears more promising for improving generation in the parallel decoding regime than for improving the aggregate downstream benchmark score, a hypothesis the final checkpoint sweep will test. Completing the paired AR evaluation will determine not merely whether these models can diffuse, but how much of their original capability survives conversion and whether that modest incremental cost is justified relative to purpose-trained diffusion models.

\section*{Reproducibility Statement}

The main text specifies the objective, trainable modules, data volume, optimization, generation, benchmarks, and scoring. The supplementary material will include resolved configurations, parameter audits, immutable code and dataset revisions, checkpoint identifiers, loss logs, generated samples, benchmark manifests, frontier-ablation outputs, and uncertainty calculations. Before submission we will complete the remaining AR cells, archive the historical dataset manifest, report conversion GPU-hours and software versions, and replace every red \texttt{XXX} marker with a measured value or an explicit limitation.

\section*{Ethics Statement}

This study uses existing language models and public benchmark or instruction datasets and does not involve human subjects. Converted models inherit factual errors, biases, unsafe behaviors, licensing constraints, and memorization risks from their parent checkpoints and training data. We report aggregate task performance rather than claiming safety or deployment readiness; release artifacts should preserve upstream licenses and document the data provenance and intended research use.

\section*{AI Use Statement}

Generative AI tools were used to assist software development and debugging, design and review of experiments, creation of analysis scripts, interpretation of preliminary results, literature discovery and summarization, and drafting and editing this manuscript. The authors manually reviewed AI-assisted code and prose, checked numerical claims against saved experimental artifacts, and verified literature claims and citations against primary sources. No benchmark value was entered solely from an unsupported model response. The authors take responsibility for the final content, claims, code, and artifacts.

\bibliography{iclr2027_conference}
\bibliographystyle{iclr2027_conference}

\appendix
\section{Result provenance and placeholder policy}

Table~\ref{tab:accuracy}, Table~\ref{tab:ppl}, and Table~\ref{tab:progression} were initially transcribed from \texttt{results/Results.xlsx} on September 18, 2026. The open-ended and merge values are archived in \texttt{Results\_open-ended-128t-32i-merged\_20260918.xlsx}. Red-font workbook entries were interpreted as 250-example estimates and black-font entries as 50-example estimates, following the workbook legend. All \texttt{XXX} markers denote information not established by the available artifact; they must be resolved rather than silently inferred before submission.

\section{Selected run identifiers}

The four 25k runs are \texttt{gemma-2-9b-mask}, \texttt{llama-3.1-8b-mask}, \texttt{qwen-2.5-7b-mask}, and \texttt{ministral-8b-mask} (the workbook shortens the last name to \texttt{ministral-mask}). The second-stage Llama run is \texttt{llama-3.1-8b-mask-continued}. The saved configurations identify the parent checkpoints listed in Table~\ref{tab:models} and record 25,000 updates for each directory.

The matched open-ended adapter results come from benchmark run \texttt{20260918T112818.173051Z--colab-validation}; the merged Llama results come from \texttt{20260918T123901.774999Z--colab-validation}. Both use 30 prompts, 128 output tokens, 32 denoising steps, seed 1234, temperature 0.7, the LLaDA-style sampler, and Microsoft Phi-4 scoring.

\section{Adapter-merging ablation}
\label{app:merge}

\begin{table}[h]
\caption{Llama adapter/merge ablation at 128 generated tokens and 32 denoising steps. PPL is token-weighted; Distinct-$n$ is averaged over answers.}
\centering
\small
\begin{tabular}{lrrrrr}
\toprule
Representation & Parameters & PPL$\downarrow$ & D-1$\uparrow$ & D-2$\uparrow$ & D-3$\uparrow$ \\
\midrule
Adapter loaded & 8.466B & 19.381 & .495 & .788 & .899 \\
Merged standalone & 8.030B & 19.283 & .498 & .807 & .918 \\
\bottomrule
\end{tabular}
\end{table}

\section{Frontier-masking ablation}
\label{app:frontier}

The selected main runs use the simpler IID corruption process: after sampling the example-level noise fraction $t$, each eligible answer or EOS-padding position is independently masked with probability $t$. In a later ablation, we tested \emph{soft-frontier masking}, intended to resemble partially completed left-to-right text while retaining stochastic exceptions. For a frontier example, masking probability increases sigmoidally with answer position around a frontier determined by $t$; the tested configuration used a 0.75 probability of selecting this branch per example, $\epsilon=0.03$, and $\tau=3.0$, with the remaining examples using IID corruption. Token losses use the corresponding inverse masking-probability correction.

The frontier variant achieved lower external Phi-4 generative perplexity than the matched IID variant, but substantially lower sliding model-token Distinct-1, Distinct-2, and Distinct-3, accompanied by visibly more immediate token and phrase repetition. This disagreement is precisely why the main experiments retain IID corruption and why generative perplexity is not used as the sole selection criterion. \xxx{Insert the matched run identifiers, checkpoint, generation NFE, token-weighted perplexity, Distinct-1/2/3 values, and representative blinded samples after consolidating the ablation artifacts.}

\end{document}
